\subsection{A five-state P5 leaf after coning the other corner}

\begin{lemma}[Proved: small monoids mapping to P5]
\label{CP2:A-small-map-to-P5}
Let $L$ be a finite monoid of order at most eight and
$\phi:L\to P_5$ a unital monoid homomorphism. Either
$B\phi$ is nullhomotopic, or $\phi$ has a monoid section.
\end{lemma}
\begin{proof}
Write $P_5=\{1\}\amalg\{e_{ij}:i,j\in\{0,1\}\}$.
A nonempty subsemigroup of its rectangular band is a full
rectangle on its occurring row and column labels: multiplying
an element of a given row by one of a given column produces
their intersection. Hence every proper unital submonoid of
$P_5$ is either trivial or the unitization of a band with one
row or one column. It has a one-sided zero and contractible
classifying space. If $\phi$ is not onto, factoring through
this image proves the first alternative.

Suppose $\phi$ is onto. Its minimal ideal $K$ maps onto the
four-element minimal band of $P_5$. To see this directly,
$\phi K$ is a nonempty ideal and contains that band. For
an element $x\in K$ mapping into the band, $LxL=K$, and
surjectivity gives $\phi K=P_5\phi(x)P_5$, exactly the band.
The row and column classes of $K$ map to row and column
classes of this band, so both its row and column counts are
at least two. If its local group $G$ were nontrivial, then
$|K|\geq2\cdot2\cdot2=8$. The global identity lies outside
$K$, because its image lies outside the target band. This
would force $|L|\geq9$. Therefore $G=1$ and $K$ itself is a
rectangular band.

The induced band homomorphism sends its row labels onto
$\{0,1\}$ and its column labels onto $\{0,1\}$. Choose one
preimage of each of these two rows and columns. Their four
intersections, together with $1_L$, form a submonoid
isomorphic to $P_5$ on which $\phi$ is the inverse
isomorphism. Its inclusion supplies the section.
\end{proof}

\begin{proposition}[Proved: the coned P5 obstruction on at most five states]
\label{CP2:A-coned-P5-five-small-opposite}
Let $S$ be a nonempty left-$L$, right-$P_5$ biset with
$|L|\leq8$ and $|S|\leq5$. Put
\[
 E=B(L,S,P_5),\qquad
 Y=\operatorname{hocofib}(E\longrightarrow BP_5),
\]
using the actual projection of the two-sided action bar.
If $\dim_{\mathbb Q}H^2(Y;\mathbb Q)\leq1$, every
rational degree-two square on $Y$ is zero.
\end{proposition}
\begin{proof}
We use first a general observation about this cofiber. If
its projection induces an injection
$H^2(BP_5;\mathbb Q)\to H^2(E;\mathbb Q)$, then every
second class of $Y$ restricts to zero on $BP_5$. It therefore
lifts to the relative group for $(Y,BP_5)$. The quotient by
$BP_5$ is a suspension, so products of these relative
positive-degree classes vanish. The same conclusion holds
when the projection is nullhomotopic: the cofiber is then
$BP_5\vee\Sigma E$, whose second products are zero since
$BP_5\simeq S^2$.

If a state is fixed by the whole right $P_5$ action, its
inclusion in the right-action bar gives a section of
$E\to BP_5$: take degree zero in the $L$ bar and use that
fixed state in every right-action simplex. No invariance
under the left $L$ action is needed for this inclusion.
The cohomology injection of the observation excludes this
case.

Suppose there is no fixed state. The band core consists of
one or two regular two-point column orbits. If it has one
orbit, both row actors agree on every state. If it has two
orbits and at most four states, the entire set is core.
For five states there is one transient state. If its two
row images lie in the same core orbit, the row actors again
agree. If they lie in different core orbits, the act is
the regular right $P_5$ act. This exhausts the possibilities.

Handle first every case with equal row actors. Let $U$ be
the full core, $E_U=B(L,U,P_5)$ and
$Y_U=\operatorname{hocofib}(E_U\to BP_5)$. The labelled
core points have the form $(b,j)$. The arrows
$h_{(b,j)}=e_{0j}$ obey
\[
 h_s m=h_{sm},\qquad h_{\ell s}=h_s
 \quad(m\in P_5,\ \ell\in L).
\]
They give the equivariant natural nullhomotopy of the
projection $E_U\to BP_5$. Thus
$Y_U\simeq BP_5\vee\Sigma E_U$, and all its second
products vanish. This argument includes two full core
orbits; no assumption on the rational homology of $BL$
has entered.

The row interchange of $P_5$ fixes every state, and hence
defines an involution $F$ of the pair $(Y,Y_U)$.
The quotient permutation module
$T=\mathbb Q[S]/\mathbb Q[U]$ is annihilated by the band.
The length-two projective $P_5$ resolution consequently
computes the relative right-action bar as $T[1]$, with
the semilinear row interchange acting by $-1$, exactly as
in Proposition~\ref{CP2:A-P5-row-independent}. Taking
derived $L$ coinvariants preserves this equivariant
comparison. Since the relative cofiber quotient is
$\Sigma(E/E_U)$, all rational relative homology
$H_*(Y,Y_U;\mathbb Q)$ has sign $-1$.

The positive eigenspace of $H_4(Y;\mathbb Q)$ is therefore
in the image of $H_4(Y_U;\mathbb Q)$. A generator of
one-dimensional $H^2(Y;\mathbb Q)$ is an eigenvector of
$F^*$, and its square has sign $+1$. It pairs to zero with
the negative fourth-homology eigenspace by its sign, and
with the positive eigenspace by the zero products on
$Y_U$. The square is zero. The case $H^2(Y)=0$ is immediate.

Finally let $S=P_5$ be the regular right act. Its commuting
left action has the form $\ell s=\phi(\ell)s$ for a
unital homomorphism $\phi:L\to P_5$. The right-regular
bar contraction identifies
$B(L,P_5,P_5)\to BL$ as a homotopy equivalence, and its
other projection identifies, up to homotopy, with $B\phi$.
This is the natural augmented regular-bar comparison: the
initial state $1$ supplies the extra degeneracy in each
right-bar fiber, while moving it across a left arrow
$\ell$ replaces its label by $\phi(\ell)$. The resulting
map on the remaining left bar is exactly $B\phi$.
By Lemma~\ref{CP2:A-small-map-to-P5}, that map is either
nullhomotopic or has a section. Both alternatives are
covered by the initial cofiber observation.
\end{proof}

\begin{corollary}[Proved: a coned small center with a five-element leaf]
\label{CP2:A-star-coned-five-leaf}
A source star with center of order at most eight, a
trivial leaf on a singleton stalk, and another leaf of
order at most five on at most five states cannot have the
rational homology and cohomology ring of $\CP^2$.
\end{corollary}
\begin{proof}
Coning the center by the singleton trivial leaf identifies
the star space with the cofiber of its other bridge
projection to the other leaf corner. If that corner is
not $P_5$, its rational second cohomology is zero by
Corollary~\ref{CP2:A-five-types}. All second classes of
the cofiber are then relative to that corner, so their
products vanish on the suspension quotient. If the corner
is $P_5$, take opposition to give it the right-action
convention of the proposition. The opposite center still
has order at most eight, and the proposition applies.
\end{proof}
